\documentclass[11pt]{beamer}
\usetheme{metropolis}

\usepackage[utf8]{inputenc}
\usepackage[T1]{fontenc}
\usepackage{amsmath,amssymb}
\usepackage{stmaryrd}
\usepackage{booktabs}
\usepackage{graphicx}
\usepackage{tikz}
\usetikzlibrary{positioning,arrows.meta,calc,decorations.pathreplacing}
\usepackage{pgfplots}
\pgfplotsset{compat=newest}
\renewcommand{\footnoterule}{}

% --- macros from the paper ---
\newcommand{\pure}[1]{\llbracket{#1}\rrbracket}
\newcommand{\vars}{\textsf{vars}}
\newcommand{\xs}{\bar{x}}
\newcommand\napprox{\not\approx}
\newcommand{\divc}[2]{#1\,\mathit{div}\,#2}
\newcommand{\modc}[2]{#1\,\mathit{mod}\,#2}
\newcommand{\solver}{\textsf{qfn2l}}
\newcommand{\CC}[3]{\mathit{C}(#1,#2,#3)}

\title{Revisiting Incremental Linearization for\\Nonlinear Integer Arithmetic}
\author{\underline{Marek Dančo} \and Karel Chvalovský \and Mikoláš Janota}
\institute{Czech Technical University in Prague}
\date{}

\begin{document}

\maketitle

\begin{frame}{Why care? Verifying smart contracts}
	A \alert{smart contract} is a program on a blockchain that itself holds
	and moves cryptocurrency.
	\begin{block}{BeautyChain (BEC) token hack, 2018}
		\small
		\begin{semiverbatim}
			\alert{uint256} amount = uint256(cnt) * \_value;  \alert{// attacker's pick}
			require(balances[msg.sender] >= amount);
		\end{semiverbatim}
		\normalsize
		\[
			\texttt{cnt}=2,\ \texttt{\_value}=2^{255}
			\Rightarrow \alert{\text{\texttt{amount}}=0}
			\Rightarrow \text{free tokens!}
		\]
	\end{block}
	\begin{itemize}
		\item This is a nonlinear arithmetic question!
	\end{itemize}
\end{frame}

\begin{frame}{Quantifier-Free Nonlinear Integer Arithmetic (NIA)}
	\begin{itemize}
		\item Quantifier-free formulas over the \alert{integers} with
		      $\{+,-,\cdot,\mathit{div},\mathit{mod},\mathit{abs}\}$
		\item Truly nonlinear:
		      \begin{itemize}
			      \item variables multiplied together ($x\cdot y$, $x^2$, $x^ky^l$)
			      \item $\mathit{div}$/$\mathit{mod}$ with a \alert{non-constant divisor}
			            ($\divc{x}{y}$, $\modc{x}{y}$)
		      \end{itemize}
		\item Satisfiability is \alert{undecidable}
		      (Hilbert's 10th problem, Matiyasevich)
	\end{itemize}
	\begin{block}{This talk}
		An incomplete but practical decision procedure for NIA
	\end{block}
\end{frame}

\begin{frame}{Simple-looking equations defeat state-of-the-art solvers}
	\[
		x^3 + y^3 + z^3 = 79
	\]
	\begin{itemize}
		\item A ``sum of three cubes'' problem
		\item A job for \alert{SMT solvers} $\leftarrow$ SAT solving + First order theories
		\item cvc5, MathSAT, Yices~2 and Z3 all fail
	\end{itemize}
\end{frame}

\begin{frame}{Can we get away with just LIA?}
	\begin{itemize}
		\item \alert{LIA} (linear arithmetic)
		      \begin{itemize}
			      \item easy for SMT solvers
			      \item Simplex algorithm + branch-and-bound
		      \end{itemize}
		\item \alert{NIA} (nonlinear, current problem)
		      \begin{itemize}
			      \item even small problems don't get solved
		      \end{itemize}
	\end{itemize}
	\begin{block}{Idea}
		Can we approximate NIA problems using LIA?
	\end{block}
\end{frame}

\begin{frame}{The idea, on an example}
	\[
		\varphi:\ \alert{x^2}\leq x \land x\geq 2
	\]
	\pause
	Replace $x^2$ by a fresh constant $p$ (a \alert{pure}).
	We get a LIA formula:
	\[
		\widehat\varphi:\ \alert{p}\leq x \land x\geq 2
	\]
	Add a refinement axiom:
	\[
		\widehat\varphi \;\land\; (x\geq 2\Rightarrow \alert{p}\geq 2x)
	\]
	\pause
	This forces:
	\[
		2x \leq \alert{p} \leq x \;\models\; x\leq 0
	\]
	unsat
\end{frame}

\begin{frame}{Our approach in one slide}
	\begin{center}
		\resizebox{\linewidth}{!}{%
			\begin{tikzpicture}[
				node distance=0.55cm and 1.15cm,
				box/.style={draw=mDarkTeal, thick, rounded corners, align=center,
						inner sep=5pt, font=\normalsize, minimum height=0.8cm,
						minimum width=1.4cm},
				res/.style={align=center, font=\normalsize\bfseries, text=mLightBrown},
				arr/.style={-{Stealth}, thick, mDarkTeal},
				lbl/.style={font=\small, midway}
				]
				\node[box] (phi) {$\varphi$};
				\node[box, right=of phi] (abs) {$\widehat\varphi\land\bigwedge\mathcal{A}$};
				\node[box, right=of abs] (lia) {LIA\\solver};
				\node[box, right=of lia] (chk) {NIA\\check};
				\node[res, right=of chk] (sat) {SAT};
				\node[res, above=1.0cm of lia] (unsat) {UNSAT};
				\node[box, draw=mLightBrown, very thick, fill=mLightBrown!12,
					font=\normalsize\bfseries, below=1.0cm of chk] (ax)
				{add axioms\\$\mathcal{A}\gets\mathcal{A}\cup\mathcal{N}$};
				\draw[arr] (phi) -- node[lbl, above] {purify} (abs);
				\draw[arr] (abs) -- (lia);
				\draw[arr] (lia) -- node[lbl, above] {sat} (chk);
				\draw[arr] (chk) -- node[lbl, above] {ok} (sat);
				\draw[arr] (lia) -- node[lbl, right] {unsat} (unsat);
				\draw[arr, mLightBrown] (chk) -- node[lbl, right] {spurious} (ax);
				\draw[arr, mLightBrown] (ax) -| ($(abs)+(0,-1.6cm)$) -- (abs);
			\end{tikzpicture}%
		}
	\end{center}
	\begin{itemize}
		\item $\widehat\varphi$ is a LIA \alert{overapproximation} of $\varphi$
		      $\Rightarrow$ \textsc{unsat} transfers
		\item A model of $\widehat\varphi$ may be \alert{spurious}
		      $\Rightarrow$ refine and repeat
	\end{itemize}
\end{frame}

\section{Algorithm}

\begin{frame}{Normalization and purification}
	\begin{block}{Normalization}
		Every nonlinear monomial is rewritten into the form $c\cdot x^ky^l$
		using fresh defined variables:\quad
		$xyz \rightsquigarrow d\,z$ with $d\approx xy$
	\end{block}
	\begin{block}{Purification}
		Every nonlinear subterm $t$ is replaced by a fresh constant $\pure{t}$:
		monomials $x^ky^l$ with $k+l\geq 2$, and $\divc{x}{y}$, $\modc{x}{y}$
		with a non-constant divisor
	\end{block}
	\begin{itemize}
		\item each monomial as a \alert{whole}: $xy^2$ gives only $\pure{xy^2}$,
		      no pures for $xy$ or $y^2$
		\item equal subterms share a pure
		\item result: a QF\_LIA formula $\widehat\varphi$ over the original
		      variables and the pures
	\end{itemize}
\end{frame}

\begin{frame}{Main loop}
	\begin{block}{\textsc{Qf-Nia-Solve}($\varphi$)}
		\begin{enumerate}
			\item $\widehat\varphi, P \gets \textsc{Purify}(\varphi)$;\quad
			      $\mathcal{A}\gets\emptyset$
			\item \textbf{loop}
			      \begin{itemize}
				      \item $\mathit{res}\gets\textsc{Check-Lia}(\widehat\varphi\land\bigwedge\mathcal{A})$
				      \item \textbf{if} $\mathit{res}=\mathit{unsat}$ \textbf{then return}
				            \alert{\textsc{unsat}}
				      \item $\mu\gets\textsc{Model}(\mathit{res})$
				      \item $\mathcal{N}\gets\textsc{Check-Nia}(\varphi,\widehat\varphi,P,\mu,\mathcal{A})$
				      \item \textbf{if} $\mathcal{N}=\emptyset$ \textbf{then return}
				            \alert{\textsc{sat}}, $\mu$
				      \item $\mathcal{A}\gets\mathcal{A}\cup\mathcal{N}$
			      \end{itemize}
		\end{enumerate}
	\end{block}
	\begin{itemize}
		\item $\mathcal{A}$ --- linear axioms accumulated across iterations
		\item Axioms are \alert{sound}: every NIA model satisfies them, so
		      \textsc{unsat} of the abstraction is conclusive
	\end{itemize}
\end{frame}

\begin{frame}{Which terms to blame? Implicant-based targeting}
	\begin{block}{\textsc{Check-Nia}($\varphi,\widehat\varphi,P,\mu,\mathcal{A}$)}
		\small
		\begin{enumerate}
			\item \textbf{if} $\mu\models_{\mathrm{NIA}}\varphi$ \textbf{then return}
			      $\emptyset$ \hfill\emph{(values of pures irrelevant)}
			\item $I\gets\textsc{Implicant}(\widehat\varphi,\mu)$
			\item $R\gets\{\pure{t}\in I \mid \mu(\pure{t})\neq\mu(t)
				      \land \pure{t}\text{ in a falsified literal}\}$
			\item $\mathit{res},\mathcal{N}\gets\textsc{Model-Fix}(\mu,R)$;\quad
			      \textbf{if} $\mathit{res}$ \textbf{then return} $\emptyset$
			\item $\mathcal{N}\gets\mathcal{N}\cup{}$congruence axioms for $R$
			\item \textbf{for each} $p\in R$:\quad
			      $\mathcal{N}\gets\mathcal{N}\cup\textsc{Axioms}(p,\mu)$
		\end{enumerate}
	\end{block}
	\vspace{-0.5em}
	\begin{itemize}
		\item \alert{Implicant}: all conjuncts, one $\mu$-satisfied disjunct per
		      disjunction
		\item only the pures in literals \alert{failing} under NIA semantics are
		      blamed --- far more targeted than prior work
	\end{itemize}
\end{frame}

\begin{frame}{Model repair}
	Before going back to the LIA solver, try to \alert{patch} the spurious model.
	\begin{block}{\textsc{Model-Fix}}
		\begin{itemize}
			\item pin the variables irrelevant to the failing literals
			\item run a few sub-iterations on the restricted LIA problem,
			      accumulating axioms
			\item if a NIA-valid model appears $\Rightarrow$ answer \textsc{sat}
			      without an outer iteration
		\end{itemize}
	\end{block}
	\begin{itemize}
		\item when it fails, the axioms gathered on the way are \alert{kept}
		\item Cimatti et al.\ return to the outer solver right after generating
		      lemmas
	\end{itemize}
\end{frame}

\begin{frame}{Example run}
	\[
		\varphi:\quad \alert<1>{x^2}\leq x \;\land\; x\geq 2
	\]
	\begin{enumerate}
		\small
		\item<1-> purify: $p:=\pure{x^2}$,\quad
		      $\widehat\varphi\gets p\leq x\land x\geq 2$
		\item<2-> LIA solve: model $\{x\mapsto 2,\ p\mapsto 1\}$
		\item<3-> NIA check: $2^2=4\neq 1$ --- \alert{spurious};
		      implicant $\{p\leq x,\ x\geq 2\}$, the literal $p\leq x$ reads
		      $4\leq 2$ \alert{false} $\Rightarrow R=\{p\}$
		\item<4-> axioms for $p$ at $v=2$, among them \textit{Lin.\,LB}:
		      \[
			      x\geq 2 \;\Rightarrow\; p\geq 2x
		      \]
		\item<5-> LIA solve: $2x\leq p\leq x$ gives $x\leq 0$, contradicting
		      $x\geq 2$ --- \alert{unsat}
		\item<6-> the abstraction is unsat $\Rightarrow$ $\varphi$ is
		      \alert{\textsc{unsat}}, after two iterations
	\end{enumerate}
\end{frame}

\section{Refinement Axioms}

\begin{frame}{Quick reminder}
	\begin{itemize}
		\item $\widehat\varphi\land\bigwedge\mathcal{A}$ is \textbf{sat}, but for
		      some nonlinear term $t$
		      \[
			      \mu(\pure{t})\;\neq\;t[\mu]
		      \]
		\item the model is \alert{spurious} --- it satisfies $\widehat\varphi$,
		      not $\varphi$
		\item Goal: linear axioms that make $\pure{t}$ track $t$ more closely and
		      \alert{rule out} this bad model
		\item Quality of the axioms decides everything:
		      \emph{tight} axioms cut iterations, \emph{cheap} axioms keep each
		      iteration fast
	\end{itemize}
\end{frame}

\begin{frame}{Axioms for $\pure{x^k}$}
	Generated when $\mu(x)=v$ but $\mu(\pure{x^k})\neq v^k$
	\begin{center}
		\footnotesize
		\begin{tabular}{ll}
			\toprule
			\textit{Sign}     & $\pure{x^k}\geq 0$ \ ($k$ even), \quad
			$x>0\iff\pure{x^k}>0$ \ ($k$ odd)                                                   \\
			\textit{Eq-zero}  & $(\pure{x^k}=0)\iff(x=0)$                                       \\
			\textit{Eq}       & $(\pure{x^k}=v^k)\iff(x=v)$                                     \\
			\textit{Gap}      & $\pure{x^k}\leq v^k \lor \pure{x^k}\geq(v+1)^k$                 \\
			\textit{Mod}      & $x\equiv v\!\!\pmod m \implies \pure{x^k}\equiv v^k\!\!\pmod m$ \\
			\midrule
			\textit{Lin.\,LB} & $x\geq v \implies \pure{x^k}\geq v^{k-1}x$                      \\
			\textit{Lin.\,UB} & $0\leq x\leq v \implies \pure{x^k}\leq v^{k-1}x$                \\
			\bottomrule
		\end{tabular}
	\end{center}
	\begin{itemize}
		\item shown for $v\geq 0$; symmetric variants for $v<0$ by the parity of $k$
		\item \textit{Sign} is model-independent, added once
		\item the \alert{secant bounds} are the new ingredient
	\end{itemize}
\end{frame}

\begin{frame}{Bounding powers by secant lines}
	\begin{columns}[c]
		\column{0.55\textwidth}
		\begin{center}
			\begin{tikzpicture}
				\begin{axis}[
						width=1.1\textwidth,
						height=1.2\textwidth,
						domain=0:4,
						xmax=4.5,
						ymax=18,
						xlabel={$x$},
						ylabel={$x^2$},
						axis lines=middle,
						xtick={1,2,3,4},
						ytick={1,4,9,16},
					]
					\addplot[mDarkTeal, very thick] {x^2};
					\addplot[blue, thick, domain=3:4] {4*x};
					\addplot[blue, densely dashed, domain=0:3] {4*x};
					\addplot[blue, thick, domain=2:4] {3*x};
					\addplot[blue, densely dashed, domain=0:2] {3*x};
					\addplot[blue, densely dashed, domain=3:4] {2*x};
					\addplot[blue, thick, domain=1:3] {2*x};
					\addplot[blue, densely dashed, domain=0:1] {2*x};
					\addplot[blue, thick, domain=0:2] {x};
					\addplot[blue, densely dashed, domain=2:4] {x};
					\addplot[red, thick] coordinates {(1,0) (1,2)};
					\addplot[red, thick] coordinates {(2,2) (2,6)};
					\addplot[red, thick] coordinates {(3,6) (3,12)};
				\end{axis}
			\end{tikzpicture}
		\end{center}
		\column{0.45\textwidth}
		The line $v^{k-1}x$ passes through the origin and $(v,v^k)$, and bounds
		$x^k$ on each side of $v$:
		\[
			x\geq v\Rightarrow \pure{x^k}\geq v^{k-1}x
		\]
		\[
			0\leq x\leq v\Rightarrow \pure{x^k}\leq v^{k-1}x
		\]
		\begin{itemize}
			\item tangent planes give only \alert{lower} bounds for convex $x^k$
			\item secants give \alert{two-sided} constraints
		\end{itemize}
	\end{columns}
\end{frame}

\begin{frame}{Bounding mixed powers $\pure{x^ky^l}$}
	Linearize $x^k$ by its secant, treat $y^l$ as a constant coefficient:
	\[
		x^ky^l \;\geq\; x^kw^l \;\geq\; v^{k-1}x\,w^l
	\]
	\begin{block}{Secant conditions $\CC{x^k}{v^k}{\downarrow/\uparrow}$}
		\centering\footnotesize
		\begin{tabular}{lll}
			\toprule
			                & $\CC{x^k}{v^k}{\downarrow}$ & $\CC{x^k}{v^k}{\uparrow}$ \\
			\midrule
			$v\geq 0$       & $x\geq v$                   & $0\leq x\leq v$           \\
			$v<0$, $k$ even & $x\leq v$                   & $v\leq x\leq 0$           \\
			$v<0$, $k$ odd  & $v\leq x\leq 0$             & $x\leq v$                 \\
			\bottomrule
		\end{tabular}
	\end{block}
	\begin{itemize}
		\item the arrows are chosen by the signs of $v^k$ and $w^l$ so that the
		      inequality survives multiplication (four cases, all proved)
		\item e.g.\ $v,w>0$:\quad
		      $x\geq v\land y\geq w \implies \pure{x^3y^2}\geq v^2w^2\,x$
		\item symmetric construction: fix $x^k$, linearize $y^l$
	\end{itemize}
\end{frame}

\begin{frame}{Tangent planes for $x\cdot y$}
	\begin{columns}[c]
		\column{0.46\textwidth}
		\begin{center}
			\includegraphics[width=\textwidth]{figs/tangent_plane.pdf}
		\end{center}
		\column{0.54\textwidth}
		\footnotesize
		Tangent plane of $xy$ at $(v,w)=(10,10)$ is
		$T = wx+vy-vw = 10x+10y-100$, and the sign of $xy-T=(x-v)(y-w)$ picks
		the side:
		\begin{align*}
			x>v\land y>w & \Rightarrow \pure{xy}>T \\
			x>v\land y<w & \Rightarrow \pure{xy}<T
		\end{align*}
		\begin{itemize}
			\item only for the bilinear case $k=l=1$
			\item = McCormick envelope at the model point
		\end{itemize}
	\end{columns}
	\vspace{0.3em}
	{\footnotesize\alert{Frontier strategy}: a per-pure bounding box of the model
		values seen so far; a point extending the box triggers tangent planes at the
		two mixed corners, so both bounds exist in every region.}
\end{frame}

\begin{frame}{Integer division and modulo}
	\small
	Axioms for $\pure{\divc{x}{y}}$ and $\pure{\modc{x}{y}}$ when $\mu(x)=v$,
	$\mu(y)=w$.
	\vspace{-0.4em}
	\begin{block}{Fix-divisor --- pinning $y$ makes it linear in $x$}
		\vspace{-0.6em}
		\begin{align*}
			y=w & \implies \pure{\divc{x}{y}}=\divc{x}{w} \\
			y=w & \implies \pure{\modc{x}{y}}=\modc{x}{w}
		\end{align*}
		\vspace{-1.2em}

		the right-hand sides are LIA, since $w$ is a numeral
	\end{block}
	\begin{block}{Large-divisor --- a big $|y|$ determines the result}
		for $v\geq 0$:
		\vspace{-0.6em}
		\begin{align*}
			x=v\land|y|>v & \implies \pure{\divc{x}{y}}=0 \\
			x=v\land|y|>v & \implies \pure{\modc{x}{y}}=v
		\end{align*}
		\vspace{-1.2em}

		for $v<0$: conclusions follow SMT-LIB, where $\mathit{mod}$ is non-negative
		($\pure{\divc{x}{y}}=\mp1$, $\pure{\modc{x}{y}}=|y|-|v|$)
	\end{block}
\end{frame}

\section{Comparison with Prior Work}

\begin{frame}{Versus Cimatti et al.\ (MathSAT5)}
	Same CEGAR skeleton --- the differences:
	\begin{center}
		\resizebox{\linewidth}{!}{%
			\footnotesize
			\begin{tabular}{lll}
				\toprule
				                 & \textbf{Cimatti et al.}           & \textbf{\solver}                    \\
				\midrule
				abstract domain  & QF\_UFLIA, $f_\times(x,y)$        & QF\_LIA, pures $\pure{t}$           \\
				congruence       & free from the UF solver           & explicit, lazy axioms               \\
				axiom generation & three rounds, stop at first hit   & all violated axioms in one pass     \\
				failing terms    & all products in falsified constr. & pures in failing implicant literals \\
				model repair     & none                              & \textsc{Model-Fix} sub-loop         \\
				axiom set        & sign, zero, neutrality,           & sign, zero, tangent \alert{plus}    \\
				                 & proportionality, tangent          & secant bounds for $x^k$, $x^ky^l$   \\
				\bottomrule
			\end{tabular}}
	\end{center}
	\begin{itemize}
		\item MathSAT5 implements that approach, so the experimental comparison
		      with MathSAT5 doubles as a comparison with it
	\end{itemize}
\end{frame}

\section{Experiments}

\begin{frame}{Setup}
	\begin{itemize}
		\item \solver{} implemented in \alert{C++} on top of the smt-switch
		      abstraction layer, using \alert{Z3 as the LIA backend}
		\item Linked against the \emph{same} Z3~4.15.4 build used as the standalone
		      Z3 baseline $\Rightarrow$ differences come from the strategy, not the
		      engine
		\item Benchmarks: the \alert{full} SMT-LIB QF\_NIA set --- 25\,444 instances
		\item 180\,s wall clock, 40\,GB memory, AMD EPYC 7513
		\item Baselines: Z3~4.15.4, MathSAT5~5.6.16, Yices~2.6.4, cvc5~1.3.2
		\item Three configurations: base, {+}\textsc{Frontier} (tangent
		      instantiation), {+}\textsc{Congr.} (lazy congruence, ablation)
	\end{itemize}
\end{frame}

\begin{frame}{Results}
	\begin{center}
		\footnotesize
		\begin{tabular}{lrr}
			\toprule
			\textbf{Solver}                     & \textbf{Total} (25\,444) & \textbf{MathProblems} (1\,100) \\
			\midrule
			\alert{\solver{}+\textsc{Frontier}} & \alert{14\,086}          & \alert{586}                    \\
			\alert{\solver}                     & \alert{14\,011}          & \alert{585}                    \\
			\alert{\solver{}+\textsc{Congr.}}   & \alert{12\,421}          & \alert{587}                    \\
			\midrule
			Z3~4.15.4                           & 19\,205                  & 118                            \\
			MathSAT5~5.6.16                     & 16\,377                  & 171                            \\
			Yices~2.6.4                         & 15\,853                  & 119                            \\
			cvc5~1.3.2                          & 12\,996                  & 157                            \\
			\bottomrule
		\end{tabular}
	\end{center}
	\begin{itemize}
		\item \alert{MathProblems} (sums of cubes): \solver{} solves \alert{53\%},
		      best baseline 16\% --- what the secant axioms for $x^k$ buy
		\item overall competitive with cvc5, behind Z3 / MathSAT / Yices
		\item frontier: small consistent gain, free.
		      Congruence: $-1\,600$ instances
	\end{itemize}
\end{frame}

\begin{frame}{Per-family picture}
	\begin{center}
		\footnotesize
		\begin{tabular}{lrr}
			\toprule
			\textbf{Family}                & \textbf{\solver{}+\textsc{Fr.}} & \textbf{best baseline} \\
			\midrule
			ITS (17\,046)                  & 4\,203/3\,908                   & 8\,154/4\,452 (Z3)     \\
			AProVE (2\,409)                & 1\,418/579                      & 1\,642/687 (Z3)        \\
			SAT14 (1\,926)                 & 1\,568/61                       & 1\,853/72 (Z3)         \\
			\midrule
			mcm (186)                      & \alert{11/0}                    & 9/0 (cvc5)             \\
			calypto (177)                  & \alert{80/97}                   & 79/97 (Z3)             \\
			UltimateAutomizerSvcomp23 (58) & \alert{7/15}                    & 8/12 (Z3)              \\
			LCTES (2)                      & \alert{0/2}                     & 0/1 (Z3)               \\
			\bottomrule
		\end{tabular}\\[0.4em]
		\scriptsize (sat/unsat solved within 180\,s)
	\end{center}
	\begin{itemize}
		\item The gap is on \alert{satisfiable} instances of large families
		\item Those are shallow, huge (hundreds of pures), with small solutions ---
		      ideal for \alert{bit-blasting}, which competitors use;
		      we must witness sat by refinement over a LIA abstraction
	\end{itemize}
\end{frame}

\begin{frame}{Conclusions and future work}
	\begin{block}{Summary}
		\begin{itemize}
			\item QF\_NIA $\to$ QF\_LIA by purification into plain constants,
			      refined by CEGAR --- no UF, no geometry
			\item new \alert{secant-based} bounds for $x^k$ and $x^ky^l$: power
			      terms as first-class abstractions
			\item implicant-based blame assignment, model repair
			\item competitive overall, \alert{dominant} on polynomial benchmarks
		\end{itemize}
	\end{block}
	\begin{block}{Future work}
		\begin{itemize}
			\item bit-blasting preprocessing for small-model sat instances
			\item pruning axioms made redundant by stronger ones
			\item new axiom schemas for further problem classes
		\end{itemize}
	\end{block}
\end{frame}

\end{document}
